\section{Data: sources, licences and pitfalls}
\label{sec:data}

\para{Free sources.} Almost everything needed to build a forecasting system is available free of charge (Table~\ref{tab:data}). Real-time forecasts come from ECMWF, whose whole real-time catalogue has been open under CC-BY-4.0 since 1 October 2025~\cite{ecmwf2025opendata}, from NOAA, DWD, M\'et\'eo-France and MET Norway. Reanalyses (ERA5~\cite{hersbach2020era5}, CERRA) and satellite irradiance records (SARAH-3, CAMS, NSRDB, PVGIS) provide the long history; national weather services and NOAA's Integrated Surface Database~\cite{smith2011isd} provide observations; the ENTSO-E Transparency Platform, SMARD, Energy-Charts and the German transmission system operators provide generation, forecasts and prices that serve both as targets and as benchmarks.

% Table 6: free data sources. Source data with URLs: data/literature/data_sources.csv
\begin{table*}[t]
\centering\small
\caption{\textbf{Free data for weather-based energy forecasting.} The last column states whether commercial use is permitted. Full list with URLs: \texttt{data/literature/data\_sources.csv}.}
\label{tab:data}
\setlength{\tabcolsep}{4pt}
\begin{tabularx}{\textwidth}{@{}L{2.6cm}L{3.9cm}YL{3.1cm}@{}}
\toprule
Purpose & Source & Content & Commercial use \\
\midrule
Real-time forecasts & ECMWF open data & IFS and AIFS, deterministic and ensemble, 0.25$^\circ$ & yes (CC-BY-4.0) \\
 & NOAA GFS, GEFS, AIGFS & global forecasts and ensembles & yes (public domain) \\
 & DWD ICON, ICON-D2, MOSMIX & fields down to 2\,km; station point forecasts & yes (attribution) \\
 & M\'et\'eo-France, MET Norway & AROME, ARPEGE, MEPS, point API & yes (open licences) \\
 & CAMS & aerosol and dust forecasts & yes (Copernicus licence) \\
Convenience API & Open-Meteo & forecasts, past forecasts, satellite radiation, ERA5 & free tier: no \\
Past forecasts & NOAA archives on AWS and Google Cloud; TIGGE & GFS, HRRR, GEFS reforecasts; multi-centre ensembles & NOAA: yes; TIGGE: no \\
Reanalysis & ERA5, ERA5-Land, CERRA, COSMO-REA6 & hourly gap-free weather since 1940 (ERA5) & yes \\
Satellite irradiance & SARAH-3, CAMS radiation, PVGIS, NSRDB & irradiance records and PV simulation & yes (attribution) \\
Observations & DWD Climate Data Center, NOAA ISD, BSRN & station weather and reference radiation & DWD, ISD: yes; BSRN: research \\
Energy and markets & ENTSO-E, SMARD, Energy-Charts, netztransparenz.de & generation, TSO forecasts, prices, imbalance & generally yes \\
Plant data & Kelmarsh and Penmanshiel (Zenodo) & 10-minute SCADA of two UK wind farms & yes (CC-BY-4.0) \\
Resource atlases & Global Solar Atlas, Global Wind Atlas; Renewables.ninja & long-term maps; simulated time series & atlases: yes; ninja: no \\
\bottomrule
\end{tabularx}
\end{table*}


\para{Licences.} Commercial use is the dividing line. NOAA data are public domain; ECMWF, DWD and Copernicus data can be used commercially with attribution; the free tier of the popular Open-Meteo API is restricted to non-commercial use~\cite{openmeteo2026terms}, as are TIGGE and Renewables.ninja. A system built for a company should either pay for a commercial tier or run directly on the open primary sources.

\para{Training data must be as issued.} A model that will run on forecasts must be trained on forecasts, with the lead time and availability it will have in operation. Training on reanalysis or observations as inputs teaches a relationship that does not exist operationally, because real inputs are worse. Archives of past forecasts exist (NOAA's GFS, HRRR and GEFS reforecast archives, ECMWF open-data mirrors, Open-Meteo's historical-forecast and previous-run interfaces), but for many regional models the only archive is the one a user builds: DWD keeps ICON files on its open-data server for about a day.

\para{Time conventions.} Store everything in UTC and record, for every source, whether a time stamp marks the start, the end or the middle of an interval, or an instant. NWP radiation is rarely instantaneous: ECMWF reports accumulations in J/m$^2$ since the start of the run, DWD's ICON reports means since the start of the run, GFS alternates 3- and 6-hour averaging windows, and Open-Meteo returns means over the preceding hour. Lemma~\ref{lem:deacc} converts the first two into interval means; getting it wrong shifts or scales the forecast systematically, most visibly in the morning and evening ramps. Solar position for an interval mean is computed at the interval's midpoint.

\para{Plant state and sensors.} Training data must be filtered for curtailment and outages, or the model learns to under-forecast; installed capacity must be tracked through extensions and repowering; clipping and export limits should be modelled explicitly. Pyranometers drift, soil and misalign; reference cells and pyranometers respond differently to spectrum and angle; nacelle anemometers sit behind the rotor. Independent satellite irradiance and clear-sky curves are the cheapest quality checks.

\para{Evaluation.} Backtests must roll forward in time (never randomly shuffle time series), cover at least one full year where possible, report errors by lead time and weather regime, compare against persistence and climatology, state the normalisation, and test differences for significance. Monitoring continues in operation: model upgrades at the weather centres shift biases overnight.
